% ─────────────────────────────────────────────────────────────────────────────
\section{Design}
% ─────────────────────────────────────────────────────────────────────────────

A set of foundational design choices guides our approach to decentralized coordination. These high-level principles govern how the fleet explores the search area and responds to discovered victims.

\paragraph{Pheromone Repulsion}
The core of our coordination system is the use of repulsive pheromones. While standard ACO \citep{ar:Dorigo1996} uses pheromones as an attractor, we invert this signal to suit a SAR mission. In our model, high pheromone levels indicate that a UAV has recently passed over a specific grid cell. By scoring candidate waypoints inversely to their pheromone intensity, we ensure that the swarm avoids redundant coverage and is naturally pushed toward unexplored territory. These pheromones evaporate over time to represent a decaying spatial memory. This allows agents to eventually return to older areas, accounting for the possibility that a victim was missed or that environmental conditions have changed.

\paragraph{Exploration--Exploitation Duality}
The repulsive pheromone signal is balanced by an attractive signal from victim detections. This interaction creates a dynamic balance between exploration and exploitation that adapts to the mission state. In the absence of detections, pheromone repulsion dominates and forces agents to spread out into unsearched space. When a UAV identifies a victim, this exploratory behaviour is immediately overridden by a direct approach to the target. Following a successful confirmation, the agent remains in the vicinity briefly to perform a local search. This strategy allows the swarm to capitalize on the spatial structure of the environment by focusing resources where they are most likely to find more victims.

\paragraph{Victim Placement Model}
\label{sec:design_victim}
We model victim placement using a Gaussian Mixture Model (GMM) to simulate the realistic spatial clustering found in disaster sites. Structural collapses or flooding events rarely result in victims being spread uniformly across a search area. Instead, victims tend to be concentrated in high-impact zones. Using a GMM allows us to generate these clusters and test the effectiveness of our local search bias. By varying the number and spread of these clusters, we can evaluate the robustness of our coordination strategy across a range of disaster scenarios, from single-point incidents to widely dispersed multi-site emergencies.
